% ════════════════════════════════════════════════════════════════════════════
%  Formation Engine — Theory Layer:
%    §3.1 Full Energy Decomposition
%    §3.3 Pairwise Mixed Interaction Structure
%    §T.1 Interaction Graph G = (V, E)
%    §T.2 Material Kernel Vector M_k
%  Printer-ready LaTeX.  Compile with pdflatex (twice for refs).
% ════════════════════════════════════════════════════════════════════════════
\documentclass[11pt, a4paper, oneside]{article}

% ── Geometry ────────────────────────────────────────────────────────────────
\usepackage[
  top=1.0in, bottom=1.0in, left=1.15in, right=1.15in,
  headheight=14pt
]{geometry}

% ── Fonts / encoding ───────────────────────────────────────────────────────
\usepackage[T1]{fontenc}
\usepackage{lmodern}
\usepackage{microtype}

% ── Math ───────────────────────────────────────────────────────────────────
\usepackage{amsmath, amssymb, amsthm, mathtools}
\usepackage{bm}

% ── Tables ─────────────────────────────────────────────────────────────────
\usepackage{booktabs}
\usepackage{array}
\usepackage{tabularx}

% ── Lists ──────────────────────────────────────────────────────────────────
\usepackage{enumitem}
\setlist{nosep, leftmargin=1.5em}

% ── Code ───────────────────────────────────────────────────────────────────
\usepackage{listings}
\usepackage{xcolor}

\lstset{
  basicstyle=\ttfamily\small,
  keywordstyle=\color{blue},
  commentstyle=\color{gray},
  stringstyle=\color{red},
  breaklines=true,
  columns=fullflexible,
  keepspaces=true,
  frame=single,
  backgroundcolor=\color{gray!8},
  xleftmargin=0.5em,
  xrightmargin=0.5em,
  language=C++
}

% ── Headers / footers ─────────────────────────────────────────────────────
\usepackage{fancyhdr}
\pagestyle{fancy}
\fancyhf{}
\fancyhead[L]{\small\textit{Formation Engine Methodology}}
\fancyhead[R]{\small\textit{Theory Layer — Energy, Graph, Material Kernel}}
\fancyfoot[C]{\thepage}
\renewcommand{\headrulewidth}{0.4pt}

% ── Section formatting ─────────────────────────────────────────────────────
\usepackage{titlesec}
\titleformat{\section}
  {\Large\bfseries}
  {§\thesection}{1em}{}
\titleformat{\subsection}
  {\large\bfseries}
  {\thesubsection}{1em}{}
\titleformat{\subsubsection}
  {\normalsize\bfseries}
  {\thesubsubsection}{1em}{}

% ── Hyperlinks ─────────────────────────────────────────────────────────────
\usepackage[
  colorlinks=true, linkcolor=black, citecolor=black, urlcolor=blue
]{hyperref}

% ── Convenience macros ─────────────────────────────────────────────────────
\newcommand{\file}[1]{\texttt{#1}}
\newcommand{\type}[1]{\texttt{#1}}
\newcommand{\field}[1]{\texttt{#1}}
\newcommand{\Ang}{\text{\AA}}
\newcommand{\eV}{\text{eV}}
\newcommand{\kcal}{\text{kcal/mol}}
\DeclareMathOperator{\tr}{tr}
\DeclareMathOperator{\diag}{diag}
\DeclareMathOperator{\sgn}{sgn}
\DeclareMathOperator{\rank}{rank}
\DeclareMathOperator{\Var}{Var}

% ════════════════════════════════════════════════════════════════════════════
\title{\textbf{Formation Engine Methodology} \\[0.5em]
\Large Theory Layer — Energy Decomposition, \\
Pairwise Interaction Structure, \\
Interaction Graph, and Material Kernel}
\author{VSEPR-SIM Project}
\date{Version 1.0 — \today}

\begin{document}
\maketitle
\thispagestyle{fancy}

% ========================================================================
\begin{abstract}
This document establishes the formal theoretical infrastructure for the
VSEPR-SIM coarse-grained simulation platform.  Four interconnected
formalisms are developed:
\begin{enumerate}
  \item \textbf{Full Energy Decomposition} (§\ref{sec:energy}):
        the nine-channel decomposition of the total potential energy
        $U_{\mathrm{tot}}$ into physically interpretable, individually
        inspectable terms.
  \item \textbf{Pairwise Mixed Interaction Structure} (§\ref{sec:pairwise}):
        the per-pair, per-species decomposition
        $U_{ij}^{xy}$ for mixed bead types $x, y \in \{a, b\}$.
  \item \textbf{Interaction Graph} (§\ref{sec:graph}):
        the graph $G = (V, E)$ encoding the combinatorial backbone of
        the bead system, with adjacency, degree, and Laplacian analysis.
  \item \textbf{Material Kernel Vector} (§\ref{sec:kernel}):
        the five-component per-bead state vector
        $\mathbf{M}_k = [\phi_k,\, \psi_k,\, \chi_k,\, \omega_k,\, E_k]^\top$
        that encodes the essential material identity of each bead or domain.
\end{enumerate}
These formalisms are deterministic, explicitly inspectable, and free of
stochastic methods.  No density-functional theory (DFT) or Monte Carlo
sampling is employed.  All quantities are computed analytically from the
coarse-grained bead state, environment observables, and registry data.

\medskip
\noindent\textbf{Design principle:}
Every mapping decision, every metric, every intermediate result must be
explicitly inspectable, traceable, and deterministic.  No hidden state.
\end{abstract}

\tableofcontents
\newpage

% ========================================================================
% §3.1  FULL ENERGY DECOMPOSITION
% ========================================================================
\section{Full Energy Decomposition}
\label{sec:energy}

\subsection{Total Energy of the System}

The total potential energy of the coarse-grained bead system is written as
the sum of nine physically distinct channels:

\begin{equation}
\boxed{
U_{\mathrm{tot}}
= U_{\mathrm{bond}}
+ U_{\mathrm{angle}}
+ U_{\mathrm{tors}}
+ U_{\mathrm{vdW}}
+ U_{\mathrm{Coul}}
+ U_{\mathrm{pol}}
+ U_{\mathrm{decay}}
+ U_{\mathrm{resp}}
+ U_{\mathrm{ext}}
}
\label{eq:Utot}
\end{equation}

Each term has a precise physical origin and a well-defined mapping to the
implementation.

\subsection{Channel Definitions}

\subsubsection{Bonded Terms:
  \texorpdfstring{$U_{\mathrm{bond}}$, $U_{\mathrm{angle}}$,
  $U_{\mathrm{tors}}$}{U\_bond, U\_angle, U\_tors}}

At the atomistic level, these terms capture covalent bond stretching,
angle bending, and torsional rotation:

\begin{align}
U_{\mathrm{bond}}  &= \sum_{(i,j) \in \mathbf{B}} \tfrac{1}{2} k_b
                       \bigl(r_{ij} - r_0\bigr)^2
  \label{eq:Ubond} \\[6pt]
U_{\mathrm{angle}} &= \sum_{\mathrm{angles}} \tfrac{1}{2} k_\theta
                       \bigl(\theta_{ijk} - \theta_0\bigr)^2
  \label{eq:Uangle} \\[6pt]
U_{\mathrm{tors}}  &= \sum_{\mathrm{dihedrals}} V_n
                       \bigl[1 + \cos(n\phi_{ijkl} - \gamma)\bigr]
  \label{eq:Utors}
\end{align}

\noindent\textbf{CG absorption.}
At the coarse-grained Level~2 resolution, direct bond, angle, and
torsion terms are \emph{absorbed} into the spherical-harmonic (SH)
orientation channels and the steric kernel.  The CG mapping
$\texttt{map\_reaction\_to\_beads}()$ distributes intramolecular
structure into the bead descriptor coefficients
$\{c_{\ell m}^{(k)}\}$.  Consequently, $U_{\mathrm{bond}}$,
$U_{\mathrm{angle}}$, and $U_{\mathrm{tors}}$ appear explicitly
only at the atomistic level; their CG image is spread across the
remaining six channels.

\subsubsection{Van der Waals:
  \texorpdfstring{$U_{\mathrm{vdW}}$}{U\_vdW}}

The isotropic Lennard-Jones 12-6 potential plus the
dispersion modulation channel~[B9]:

\begin{equation}
U_{\mathrm{vdW}} = \sum_{i < j}
  4\varepsilon_{ij}\left[
    \left(\frac{\sigma_{ij}}{r_{ij}}\right)^{\!12}
  - \left(\frac{\sigma_{ij}}{r_{ij}}\right)^{\!6}
  \right]
  \cdot g_d(\bar\eta_{ij})
\label{eq:UvdW}
\end{equation}

where $g_d(\bar\eta) = 1 + \gamma_d \cdot \bar\eta$ is the dispersion
modulation factor and $\bar\eta_{ij} = \tfrac{1}{2}(\eta_i + \eta_j)$
is the mean slow state of the pair.

\subsubsection{Electrostatic:
  \texorpdfstring{$U_{\mathrm{Coul}}$}{U\_Coul}}

The Coulombic interaction with environment-responsive screening~[B8]:

\begin{equation}
U_{\mathrm{Coul}} = \sum_{i < j}
  k_e\,\frac{q_i\,q_j}{r_{ij}}
  \cdot g_e(\bar\eta_{ij})
\label{eq:UCoul}
\end{equation}

where $g_e(\bar\eta) = 1 + \gamma_e \cdot \bar\eta$ and
$\gamma_e < 0$ implements electrostatic screening in dense
environments.

\subsubsection{Polarisation:
  \texorpdfstring{$U_{\mathrm{pol}}$}{U\_pol}}

The induced-dipole energy from the QM Level-0 polarisability proxy:

\begin{equation}
U_{\mathrm{pol}} = -\frac{1}{2}\sum_i \alpha_i\,|\mathbf{E}_i^{\mathrm{local}}|^2
\label{eq:Upol}
\end{equation}

where $\alpha_i$ is the polarisability proxy (in~$\Ang^3$) from
\type{QMDescriptor}, and $\mathbf{E}_i^{\mathrm{local}}$ is the
local electric field at bead~$i$ from all neighbours.

\subsubsection{Decay Dissipation:
  \texorpdfstring{$U_{\mathrm{decay}}$}{U\_decay}}

Energy stored in the slow-state relaxation dynamics that has not yet
dissipated to steady state:

\begin{equation}
U_{\mathrm{decay}} = -\frac{1}{\tau}\sum_i
  \bigl(f_i - \eta_i\bigr)^2
\label{eq:Udecay}
\end{equation}

where $f_i$ is the target function $f(\hat\rho_i, \hat P_{2,i})$
and $\eta_i$ is the current slow state.  This term vanishes at
convergence ($f_i = \eta_i$ for all~$i$).

An additional pairwise decay-coupled term captures interfacial
dissipation from $\eta$-mismatch between neighbouring beads:

\begin{equation}
U_{\mathrm{decay}}^{\mathrm{coupled}}
= -\gamma_d \sum_{i < j}
  |\Delta\eta_{ij}|\cdot S_{\mathrm{sw}}(r_{ij};\, r_c, \delta)
\label{eq:Udecay_coupled}
\end{equation}

where $S_{\mathrm{sw}}$ is the smooth switching function and
$\Delta\eta_{ij} = \eta_i - \eta_j$.  Pairs with mismatched
$\eta$ (one ordered, one disordered) experience a dissipative
coupling that drives them toward compatible states — the CG
equivalent of interfacial energy at domain boundaries.

\subsubsection{Environment-Responsive Coupling:
  \texorpdfstring{$U_{\mathrm{resp}}$}{U\_resp}}

The excess energy arising from environment modulation — the difference
between the modulated and unmodulated interaction:

\begin{equation}
U_{\mathrm{resp}} = \sum_{i < j}\sum_k
  \bigl(g_k(\bar\eta_{ij}) - 1\bigr)\cdot U_{ij,k}^{\mathrm{bare}}
= \sum_{i < j}\sum_k
  \gamma_k\,\bar\eta_{ij}\cdot U_{ij,k}^{\mathrm{bare}}
\label{eq:Uresp}
\end{equation}

where the sum over $k$ runs over channels
$\{\mathrm{steric},\,\mathrm{elec},\,\mathrm{disp}\}$
and $U_{ij,k}^{\mathrm{bare}}$ is the unmodulated base interaction.

\subsubsection{External Field:
  \texorpdfstring{$U_{\mathrm{ext}}$}{U\_ext}}

Coupling to an externally imposed field (electric, magnetic, gravitational,
or confinement potential):

\begin{equation}
U_{\mathrm{ext}} = \sum_i V_{\mathrm{ext}}(\mathbf{R}_i)
\label{eq:Uext}
\end{equation}

This channel is a placeholder for future extensions.  Currently
$U_{\mathrm{ext}} = 0$ in all simulations.

\subsection{Energy Decomposition Diagnostic}

The \type{EnergyDecomposition} struct provides per-channel access
to all terms at every step.  The total is reconstructed additively:

\begin{equation}
U_{\mathrm{tot}} = \sum_{k=0}^{7} U_k
\label{eq:Utot_sum}
\end{equation}

where each $U_k$ is tagged with an \type{EnergyChannel} enum.
Implementation: \file{coarse\_grain/theory/energy\_decomposition.hpp}.

\subsection{Mapping to Implementation}

\begin{center}
\small
\begin{tabularx}{\textwidth}{lllX}
\toprule
\textbf{Channel} & \textbf{Tag} & \textbf{6+9 Unit} & \textbf{Source} \\
\midrule
$U_{\mathrm{vdW}}$   & \type{VdW}      & S4, B9 & LJ 12-6 + $g_d$ \\
$U_{\mathrm{Coul}}$  & \type{Coulomb}  & S4, B8 & Coulomb + $g_e$ \\
$U_{\mathrm{steric}}$& \type{Steric}   & S4, B7 & SH steric kernel + $g_s$ \\
$U_{\mathrm{orient}}$& \type{Orient}   & B4     & SH overlap integral \\
$U_{\mathrm{pol}}$   & \type{Pol}      & QM L0  & $-\tfrac{1}{2}\alpha|\mathbf{E}|^2$ \\
$U_{\mathrm{decay}}$ & \type{Decay}    & B6     & $\eta$-relaxation + coupled \\
$U_{\mathrm{resp}}$  & \type{Resp}     & B7--B9 & $(g_k - 1)\cdot U^{\mathrm{bare}}$ \\
$U_{\mathrm{ext}}$   & \type{External} & —      & Future placeholder \\
\bottomrule
\end{tabularx}
\end{center}


% ========================================================================
% §3.3  PAIRWISE MIXED INTERACTION STRUCTURE
% ========================================================================
\section{Pairwise Mixed Interaction Structure}
\label{sec:pairwise}

\subsection{Per-Pair Decomposition}

Each pairwise interaction between beads $i$ and $j$ of species
$x, y \in \{a, b\}$ is decomposed into five terms:

\begin{equation}
\boxed{
U_{ij}^{xy}
= U_{ij,\mathrm{steric}}^{xy}
+ U_{ij,\mathrm{elec}}^{xy}
+ U_{ij,\mathrm{disp}}^{xy}
+ U_{ij,\mathrm{orient}}^{xy}
+ U_{ij,\mathrm{decay\text{-}coupled}}^{xy}
}
\qquad x, y \in \{a, b\}
\label{eq:Uij_xy}
\end{equation}

\subsection{Channel Formulas}

\subsubsection{Steric Channel}

The steric interaction is the SH-expanded overlap with
environment-modulated kernel:

\begin{equation}
U_{ij,\mathrm{steric}}^{xy}
= \sum_{\ell=0}^{\ell_{\max}}
  K_s(\ell, r_{ij})\cdot
  g_s(\bar\eta_{ij})\cdot
  \sum_{m=-\ell}^{\ell} c_{\ell m}^{(s,x,i)}\,\tilde{c}_{\ell m}^{(s,y,j)}
\label{eq:Uij_steric}
\end{equation}

where $\tilde{c}$ denotes coefficients of bead~$j$ rotated into the
frame of bead~$i$ via the Wigner D-matrix, and
$g_s = 1 + \gamma_s\cdot\bar\eta_{ij}$ is the steric modulation factor.

\subsubsection{Electrostatic Channel}

\begin{equation}
U_{ij,\mathrm{elec}}^{xy}
= \sum_{\ell=0}^{\ell_{\max}}
  K_e(\ell, r_{ij})\cdot
  g_e(\bar\eta_{ij})\cdot
  \sum_{m=-\ell}^{\ell} c_{\ell m}^{(e,x,i)}\,\tilde{c}_{\ell m}^{(e,y,j)}
\label{eq:Uij_elec}
\end{equation}

The electrostatic kernel $K_e(\ell, r)$ decays as $r^{-(2\ell+1)}$
for the multipole expansion and as $r^{-1}$ for the monopole.
The modulation $g_e < 1$ in dense environments (screening).

\subsubsection{Dispersion Channel}

\begin{equation}
U_{ij,\mathrm{disp}}^{xy}
= \sum_{\ell=0}^{\ell_{\max}}
  K_d(\ell, r_{ij})\cdot
  g_d(\bar\eta_{ij})\cdot
  \sum_{m=-\ell}^{\ell} c_{\ell m}^{(d,x,i)}\,\tilde{c}_{\ell m}^{(d,y,j)}
\label{eq:Uij_disp}
\end{equation}

The dispersion kernel $K_d(\ell, r)$ has $r^{-6}$ leading order.
Enhancement under ordering: $g_d > 1$ when $\bar\eta > 0$.

\subsubsection{Orientational Channel}

\begin{equation}
U_{ij,\mathrm{orient}}^{xy}
= \lambda_{\mathrm{orient}}\cdot
  \sum_{k\in\{s,e,d\}}
  \sum_{\ell,m} c_{\ell m}^{(k,x,i)}\,\tilde{c}_{\ell m}^{(k,y,j)}
  \cdot w_\ell(r_{ij})
\label{eq:Uij_orient}
\end{equation}

This is the pure orientational overlap contribution, summing the SH
dot products across all channels with distance-dependent weights
$w_\ell(r) = S_{\mathrm{sw}}(r;\, r_c, \delta)$.

\subsubsection{Decay-Coupled Channel}

\begin{equation}
U_{ij,\mathrm{decay\text{-}coupled}}^{xy}
= -\gamma_d\cdot |\eta_i - \eta_j|\cdot S_{\mathrm{sw}}(r_{ij};\, r_c, \delta)
\label{eq:Uij_decay}
\end{equation}

This captures the dissipative penalty for $\eta$-mismatch across an
interface, driving neighbouring beads toward coherent internal states.

\subsection{Species Mixing Rules}

For pairs of different species $x \neq y$, the interaction parameters
follow the standard combining rules:

\begin{align}
\sigma_{xy} &= \tfrac{1}{2}(\sigma_x + \sigma_y)
  &\text{(Lorentz rule)} \label{eq:lorentz} \\
\varepsilon_{xy} &= \sqrt{\varepsilon_x \cdot \varepsilon_y}
  &\text{(Berthelot rule)} \label{eq:berthelot}
\end{align}

For metallic systems, the alloy proxy uses Vegard's law:

\begin{equation}
X_{\mathrm{alloy}} = (1 - x_B)\,X_A + x_B\,X_B
\label{eq:vegard}
\end{equation}

for properties $X \in \{a_0,\, M,\, E_{\mathrm{coh}},\, \mathrm{CN}_{\mathrm{bulk}}\}$.

\subsection{Modulation Invariant}

The modulation factor must preserve the sign of the base kernel:

\begin{equation}
|1 + \gamma_k \cdot \bar\eta| > 0
\qquad \forall\; k \in \{s, e, d\},\;\; \bar\eta \in [0, 1]
\label{eq:modulation_invariant}
\end{equation}

This is enforced numerically with a floor of $10^{-10}$.

\subsection{Diagnostic Record}

The \type{PairInteraction} struct carries all five energy terms,
the three modulation factors $g_s, g_e, g_d$, the mean $\bar\eta$,
the separation $r_{ij}$, and the species tags $x, y$.  Every pairwise
interaction is individually inspectable.

Implementation: \file{coarse\_grain/theory/energy\_decomposition.hpp}.


% ========================================================================
% §T.1  INTERACTION GRAPH
% ========================================================================
\section{Interaction Graph}
\label{sec:graph}

\subsection{Graph Definition}

The bead system defines an undirected interaction graph:

\begin{equation}
\boxed{G = (V, E), \qquad |V| = N}
\label{eq:graph}
\end{equation}

where $V = \{1, 2, \ldots, N\}$ is the vertex set (one vertex per bead)
and $E \subseteq \binom{V}{2}$ is the edge set, determined by spatial
proximity.

\subsection{Adjacency}

The adjacency matrix $\mathbf{A} \in \{0, 1\}^{N \times N}$ is defined by:

\begin{equation}
A_{ij} = \begin{cases}
  1 & \text{if } (i, j) \in E \\[3pt]
  0 & \text{if } (i, j) \notin E
\end{cases}
\label{eq:adjacency}
\end{equation}

An edge exists between beads $i$ and $j$ if and only if they are within
the interaction cutoff:

\begin{equation}
(i, j) \in E \;\;\Longleftrightarrow\;\;
  |\mathbf{R}_i - \mathbf{R}_j| < r_{\mathrm{cut}}
\label{eq:edge_criterion}
\end{equation}

For weighted graphs, $A_{ij}$ is replaced by the switching function:

\begin{equation}
A_{ij}^{(w)} = S_{\mathrm{sw}}(r_{ij};\, r_c, \delta)
\label{eq:weighted_adj}
\end{equation}

\subsection{Degree}

The degree of vertex $i$ is the number of edges incident to it:

\begin{equation}
\deg(i) = \sum_{j=1}^{N} A_{ij}
\label{eq:degree}
\end{equation}

This maps directly to the coordination number $C_i$ from the
environment-responsive bead state:

\begin{equation}
C_i \equiv \deg(i)
\label{eq:C_is_degree}
\end{equation}

\subsection{Degree Matrix and Graph Laplacian}

The degree matrix is the diagonal matrix of degrees:

\begin{equation}
\mathbf{D} = \diag\bigl(\deg(1),\, \deg(2),\, \ldots,\, \deg(N)\bigr)
\label{eq:degree_matrix}
\end{equation}

The combinatorial graph Laplacian is:

\begin{equation}
\boxed{\mathbf{L} = \mathbf{D} - \mathbf{A}}
\label{eq:laplacian}
\end{equation}

with entries:

\begin{equation}
L_{ij} = \begin{cases}
  \deg(i) & \text{if } i = j \\[3pt]
  -A_{ij} & \text{if } i \neq j
\end{cases}
\label{eq:laplacian_entries}
\end{equation}

\subsection{Spectral Properties}

The eigenvalues $\lambda_1 \leq \lambda_2 \leq \cdots \leq \lambda_N$
of $\mathbf{L}$ characterise the connectivity of the graph:

\begin{itemize}
\item $\lambda_1 = 0$ always (eigenvector $\mathbf{1}$, the constant vector).
\item The multiplicity of $\lambda = 0$ equals the number of connected
      components.
\item $\lambda_2 > 0$ if and only if $G$ is connected.
      $\lambda_2$ is the \textbf{algebraic connectivity} (Fiedler value).
\item $\tr(\mathbf{L}) = \sum_i \deg(i) = 2|E|$.
\end{itemize}

For the interaction graph, the Fiedler value measures how strongly the
bead system is connected.  A small $\lambda_2$ indicates a system close
to fragmentation (near-critical domain boundary).

\subsection{Graph Diagnostics}

The \type{GraphDiagnostics} struct provides:

\begin{center}
\small
\begin{tabular}{lll}
\toprule
\textbf{Quantity} & \textbf{Formula} & \textbf{Physical meaning} \\
\midrule
$|V|$           & $N$                        & Bead count \\
$|E|$           & edge count                 & Interaction pair count \\
$\langle k\rangle$ & $2|E|/N$                & Mean coordination \\
$\max(k)$       & $\max_i \deg(i)$           & Most connected bead \\
$\Var(k)$       & sample variance of degrees & Structural heterogeneity \\
$\langle r\rangle$ & mean $r_{ij}$ over $E$  & Mean interaction distance \\
$r_{\min}$      & $\min r_{ij}$              & Closest pair \\
$r_{\max}$      & $\max r_{ij}$              & Edge at cutoff boundary \\
\bottomrule
\end{tabular}
\end{center}

Implementation: \file{coarse\_grain/theory/graph\_topology.hpp}.


% ========================================================================
% §T.2  MATERIAL KERNEL VECTOR
% ========================================================================
\section{Material Kernel Vector}
\label{sec:kernel}

\subsection{Definition}

The material kernel vector $\mathbf{M}_k$ is the five-component
per-bead (or per-domain) representation that encodes the essential
material identity:

\begin{equation}
\boxed{
\mathbf{M}_k
= \begin{pmatrix}
  \phi_k \\[3pt]
  \psi_k \\[3pt]
  \chi_k \\[3pt]
  \omega_k \\[3pt]
  E_k
\end{pmatrix}
\in \mathbb{R}^5
}
\label{eq:Mk}
\end{equation}

\subsection{Component Definitions}

\subsubsection{Electrostatic Potential \texorpdfstring{$\phi_k$}{φ\_k}}

The electrostatic potential at the bead centre, from the Coulombic field
of all neighbours modulated by the environment coupling:

\begin{equation}
\phi_k = \sum_{j \neq k}
  k_e \, \frac{q_j \cdot g_e(\bar\eta_{kj})}{r_{kj}}
\qquad [\kcal/e]
\label{eq:phi}
\end{equation}

Source: \field{QMDescriptor::phi\_elec} or, in the metal-registry
fallback, the work function $\varphi$ of the element.

\subsubsection{Structural Coherence
  \texorpdfstring{$\psi_k$}{ψ\_k}}

The slow-state variable encoding the degree of structural order
attained by the bead through relaxation dynamics:

\begin{equation}
\psi_k = \eta_k \in [0, 1]
\label{eq:psi}
\end{equation}

where $\eta_k$ is governed by the first-order relaxation equation:

\begin{equation}
\frac{d\eta_k}{dt} = \frac{f_k - \eta_k}{\tau}
\label{eq:eta_ode}
\end{equation}

At convergence, $\psi_k \to f_k = \alpha\hat\rho_k + \beta\hat P_{2,k}$.

\subsubsection{Electronegativity
  \texorpdfstring{$\chi_k$}{χ\_k}}

The electronegativity of the bead on the Pauling scale, encoding
chemical identity and electron-drawing power:

\begin{equation}
\chi_k = \bar\chi_{\mathrm{atoms}} = \frac{1}{N_{\mathrm{atoms}}}
  \sum_{a \in \mathrm{bead}\;k} \chi_a^{\mathrm{Pauling}}
\label{eq:chi}
\end{equation}

Source: \field{QMDescriptor::chi\_mean} or
\field{MetalRecord::electronegativity\_pauling}.

\subsubsection{Orbital Overlap
  \texorpdfstring{$\omega_k$}{ω\_k}}

The orbital overlap proxy, measuring the degree of covalent bonding
character in the bead's local neighbourhood:

\begin{equation}
\omega_k = \frac{1}{N_{\mathrm{nbr}}}
  \sum_{j \in \mathrm{nbr}(k)}
  \exp\!\left(-\frac{r_{kj}}{r_{\mathrm{overlap}}}\right)
  \cdot \frac{\chi_k \cdot \chi_j}{\chi_k + \chi_j}
\qquad \in [0, 1]
\label{eq:omega}
\end{equation}

where $r_{\mathrm{overlap}} = \frac{1}{2}(\sigma_k + \sigma_j)$ is the
mean LJ radius.
Source: \field{QMDescriptor::omega\_overlap}.

\subsubsection{Local Energy Density
  \texorpdfstring{$E_k$}{E\_k}}

The per-bead share of the total potential energy:

\begin{equation}
E_k = \frac{U_{\mathrm{tot}}}{N}
\qquad [\kcal]
\label{eq:Ek}
\end{equation}

This is the mean-field approximation.  A pair-resolved variant assigns:

\begin{equation}
E_k^{\mathrm{pair}} = \frac{1}{2}\sum_{j \neq k}
  U_{kj}^{\mathrm{tot}}
\label{eq:Ek_pair}
\end{equation}

\subsection{Properties of \texorpdfstring{$\mathbf{M}_k$}{M\_k}}

\begin{enumerate}
\item \textbf{Deterministic}: computed entirely from CG state +
      registry data.  No stochastic sampling.
\item \textbf{Low-dimensional}: 5 real-valued components — tractable
      for comparison, clustering, and visualisation.
\item \textbf{Rotationally invariant}: all components are scalars;
      $\mathbf{M}_k$ is independent of the global coordinate frame.
\item \textbf{Basis for comparison}: the Euclidean distance
      $\|\mathbf{M}_a - \mathbf{M}_b\|_2$ measures material
      dissimilarity between two beads or domains.
\item \textbf{Basis for classification}: clustering in
      $\mathbb{R}^5$ identifies material families.
\end{enumerate}

\subsection{Distance and Similarity}

The Euclidean distance between two kernels is:

\begin{equation}
d(\mathbf{M}_a, \mathbf{M}_b)
= \left[
  (\phi_a - \phi_b)^2
+ (\psi_a - \psi_b)^2
+ (\chi_a - \chi_b)^2
+ (\omega_a - \omega_b)^2
+ (E_a - E_b)^2
\right]^{1/2}
\label{eq:dist}
\end{equation}

The cosine similarity is:

\begin{equation}
\cos\theta = \frac{\mathbf{M}_a \cdot \mathbf{M}_b}
  {\|\mathbf{M}_a\|\cdot\|\mathbf{M}_b\|}
\label{eq:cosine}
\end{equation}

The \type{MaterialDistance} struct provides both measures plus
per-component differences for diagnostic inspection.

\subsection{Domain Aggregation}

For a Level~3 domain containing $n$ member beads, the aggregate
kernel is the arithmetic mean:

\begin{equation}
\mathbf{M}_{\mathrm{domain}} = \frac{1}{n}\sum_{k=1}^{n} \mathbf{M}_k
\label{eq:Mdomain}
\end{equation}

This is the kernel vector that flows from L3 to the macro-DM layer.

\subsection{Kernel Spectrum}

The system-wide distribution of kernels is characterised by:

\begin{align}
\langle\mathbf{M}\rangle &= \frac{1}{N}\sum_{k=1}^{N} \mathbf{M}_k
  &\text{(component-wise mean)} \label{eq:Mmean} \\
\Var(\mathbf{M}) &= \frac{1}{N-1}\sum_{k=1}^{N}
  (\mathbf{M}_k - \langle\mathbf{M}\rangle)
  \odot
  (\mathbf{M}_k - \langle\mathbf{M}\rangle)
  &\text{(component-wise variance)} \label{eq:Mvar}
\end{align}

where $\odot$ denotes the Hadamard (element-wise) product.

A low variance in all components indicates a homogeneous material;
high variance in $\psi$ (structural coherence) indicates coexisting
ordered and disordered domains.

Implementation: \file{coarse\_grain/theory/material\_kernel.hpp}.


% ========================================================================
% §T.3  INTERCONNECTION MAP
% ========================================================================
\section{Interconnection of Formalisms}
\label{sec:interconnection}

The four formalisms developed in this document are not independent;
they form a connected theoretical infrastructure:

\begin{center}
\begin{tabular}{lcl}
\textbf{Formalism} & $\longrightarrow$ & \textbf{Consumer} \\
\midrule
$G = (V, E)$ & provides & edge iteration for $\sum_{i < j}$ in §\ref{sec:energy} \\
$U_{ij}^{xy}$ & provides & per-pair energy for $E_k$ in $\mathbf{M}_k$ \\
$\deg(i)$     & provides & coordination $C_i$ for environment state \\
$\eta_i$      & provides & $\psi_k$ component of $\mathbf{M}_k$ \\
$\mathbf{M}_k$ & provides & input to L3 domain aggregation \\
$\mathbf{M}_{\mathrm{domain}}$ & provides & input to macro-DM \\
\end{tabular}
\end{center}

\noindent The data flow is:

\begin{equation}
\text{Atomistic}
\;\xrightarrow{\text{CG map}}\;
G = (V,E)
\;\xrightarrow{\text{§3.3 pairs}}\;
U_{\mathrm{tot}}
\;\xrightarrow{\text{§T.2}}\;
\mathbf{M}_k
\;\xrightarrow{\text{L3 agg}}\;
\mathbf{M}_{\mathrm{domain}}
\;\xrightarrow{\text{handoff}}\;
\text{Macro-DM}
\label{eq:dataflow}
\end{equation}

Every intermediate is explicitly inspectable.  No hidden state
propagates through this pipeline.


% ========================================================================
% REFERENCES
% ========================================================================
\newpage
\section*{References}
\addcontentsline{toc}{section}{References}

\begin{enumerate}[label={[\arabic*]}, leftmargin=2em]
\item Rappé, A.~K., Casewit, C.~J., Colwell, K.~S., Goddard, W.~A.,
      Skiff, W.~M.,
      ``UFF, a Full Periodic Table Force Field for Molecular Mechanics
      and Molecular Dynamics Simulations,''
      \textit{J.~Am.\ Chem.\ Soc.}\ \textbf{114}, 10024--10035 (1992).

\item Kittel, C., \textit{Introduction to Solid State Physics},
      8th~ed.\ (Wiley, 2005).

\item CRC \textit{Handbook of Chemistry and Physics},
      105th~ed.\ (CRC Press, 2024).

\item Tyson, W.~R.\ \& Miller, W.~A.,
      ``Surface Free Energies of Solid Metals,''
      \textit{Surf.\ Sci.}\ \textbf{62}, 267--276 (1977).

\item Baletto, F.\ \& Ferrando, R.,
      ``Structural Properties of Nanoclusters,''
      \textit{Rev.\ Mod.\ Phys.}\ \textbf{77}, 371--423 (2005).

\item NIST XCOM Photon Cross Section Database,
      \url{https://physics.nist.gov/PhysRefData/Xcom/html/xcom1.html}.

\item Robinson, M.~T.\ \& Torrens, I.~M.,
      ``Computer Simulation of Atomic-Displacement Cascades in Solids,''
      \textit{Phys.\ Rev.\ B}\ \textbf{9}, 5008--5024 (1974).
      (Kinchin--Pease / NRT displacement model.)

\item Debye, P.,
      ``Zerstreuung von Röntgenstrahlen,''
      \textit{Ann.\ Phys.}\ \textbf{351}, 809--823 (1915).
      (Debye scattering equation.)

\item Fiedler, M.,
      ``Algebraic Connectivity of Graphs,''
      \textit{Czech.\ Math.\ J.}\ \textbf{23}, 298--305 (1973).
\end{enumerate}


\end{document}
